\chapter{Elementos pedagógicos y programas de los cursos}


\section{Análisis pedagógico}

En este apartado se presenta la caracterización del marco pedagógico que se considera viable para la carrera de Matemática, al tomar en cuenta cuáles aspectos se deben privilegiar por cada enfoque pedagógico y al considerar el aporte de las experiencias de los docentes que, en los últimos años, han mencionado la importancia de transformar los espacios de aprendizaje y la forma en la que se concibe la relación docente-estudiante. 

\subsection{Aspectos generales}

Los procesos de enseñanza y aprendizaje se van determinando por medio de prácticas que siguen los aspectos teóricos establecidos por la pedagogía, puesto que es una disciplina de estudio donde se “indaga, reflexiona y sistematizan las prácticas educativas y formativas, desde diferentes estatutos epistemológicos y metodológicos” \cite[p.~10]{AriasFrancis2011}. De esta disciplina, se incorporan modelos pedagógicos que permiten una aproximación de la realidad a la que se debe enfrentar una comunidad, desde un punto de vista formativo, considerando que se establezca una forma de trabajo ordenada y coherente, cuya comprensión, según \cite{AriasFrancis2011}, “implica reconocer las huellas que sirven como fundamento para la reflexión y la investigación, y como marco de referencia para interpretar y comprender implicaciones, limitaciones y alcances” (p.~19).

Para un acercamiento a una visión de lo que supone un componente pedagógico en la carrera de Matemática, a partir de las experiencias que expresaron los docentes en los talleres, y, según los marcos referenciales presentados anteriormente, se han analizado las características personales y colectivas que se han considerado en diferentes instantes al enseñar matemática, y se coincide que se requiere un camino de mejora en la forma de aplicar la pedagogía en la matemática, precisamente con la adquisición de estrategias que permitan una mejor toma de decisiones en esta disciplina.  
Considerando que el contexto educativo pertenece al siglo XXI, hay un interés por conocer el valor práctico de la matemática, inclusive más allá de la posibilidad de resolver problemas. Gil \cite{gil2005reflexiones} expresó que mostrar ese valor práctico es tarea del docente, así como el ser capaz de despertar la motivación necesaria, que lleve al estudiante a aprender más y mejor. Este es un gran reto, lo cual invita a que diversos profesionales se unan para que se logren activar competencias que conlleven el mundo de las relaciones entre la matemática, la visión del universo y la tecnología. Según el año de estudio en el que se encuentre el estudiante, el docente deberá diseñar actividades adecuadas para ese momento.  


\subsection{Enfoques pedagógicos}

Para limitar la propuesta del perfil de salida de la carrera de Matemática, se vislumbra un modelo pedagógico que incorpore varios enfoques, que estén acordes a la propuesta curricular y cuyo personal docente es el que lo va a ejecutar. Se considera primordialmente el enfoque cognitivo, pero también se destacan aspectos del conductista y del crítico, puesto que, en una etapa inicial de la carrera, se considera que el enfoque cognitivo es más complejo de implementar satisfactoriamente; y hacia la parte final de la carrera se busca que el estudiante realice cognitivamente los procesos de aprendizaje.

La coyuntura actual de la educación secundaria en nuestro país, en la que el estudiante llega a la Universidad con conocimientos matemáticos casi nulos y conceptualmente deteriorados, dificultan que la carrera comience con un enfoque netamente cognitivo o crítico. Los primeros ciclos de la carrera deberán contener una alta proporción de aprendizaje conductista, con pequeños matices de los enfoques cognitivo y crítico. 

En esta etapa temprana, el enfoque cognitivo puede irse implantando de manera paulatina mediante talleres y actividades lúdicas o de estudio independiente, incrustadas dentro de los cursos teórico-prácticos. Más que partir de los escasos conocimientos previos en matemática, el estudiante podría partir del conocimiento de la realidad, para comenzar a entenderla desde una perspectiva matemática. La solución de algunos problemas sociales concretos mediante modelos matemáticos simples puede ayudar a ir implantando el modelo crítico, el cual también se puede fortalecer paulatinamente mediante seminarios de corte histórico, la realización de foros de discusión y el aprendizaje colaborativo. A medida que se avanza en la carrera, los componentes cognitivo y crítico, pero, sobre todo, el primero, deberán irse acentuando de manera transversal. 

Esta presentación de la secuencia de la enseñanza subyacente con los enfoques requiere precisar cada uno de estos, tanto en cómo se conciben, como las premisas pedagógicas correspondientes. González \cite{Gonzalez2014} presenta un material para innovar en docencia universitaria y esta descripción de los enfoques considera apartados presentes en ese documento que fueron considerados como apropiados para orientar al docente de Matemática. 

\subsubsection*{Enfoque cognitivo}

Este enfoque plantea al conocimiento como un acto interior del desarrollo humano, para comprender la sociedad, adaptarse y actuar en ella.  Difiere del conductismo al plantear que el conocimiento no se obtiene ni se posee, sino que se construye y se interioriza, lo cual genera personas más apropiadas de su conocimiento, además de que construyen, a su vez, un criterio acerca de lo que aprenden. 

Las premisas pedagógicas en este enfoque son: 


\begin{itemize}
    
    \color{blue} \item \color{black} El docente tiene el conocimiento para enseñar.  

    
     \color{blue} \item \color{black} Los estudiantes son personas activas, con conocimientos previos y comprenden su contexto. 
     
    \color{blue} \item \color{black}  Los estudiantes construyen el conocimiento, lo asimilan y lo acomodan, relacionándolo con sus aprendizajes anteriores, de manera que se vuelve significativo. 
    
     \color{blue} \item \color{black} El conocimiento es una verdad que debe ser relacionada con los aprendizajes anteriores y contrastada con la realidad. 
     
     \end{itemize}
     
La persona que tiene a cargo la docencia, y que practica este enfoque, se plantea objetivos cuyo fin esencial es la construcción de conocimientos, de manera empírica, al relacionarlos con la realidad mediante actividades de aprendizaje en las que el docente se convierte en una guía que toma en cuenta los aprendizajes previos de sus estudiantes. Su tarea es velar porque los estudiantes enriquezcan su experiencia mediante el razonamiento, la solución de problemas, los estudios de caso, la investigación bibliográfica, las prácticas y otras actividades que vayan de lo concreto a lo abstracto, en función del grado de maduración del sujeto; revisiones periódicas a conceptos ya aprendidos; y, retroalimentación a partir de un problema ya resuelto. 

\subsubsection*{Enfoque conductista}

El conductismo dice que el conocimiento es evidenciable. Por ejemplo, aprender a escribir se evidencia con la posibilidad de hacerlo, y, por lo tanto, el conocimiento se obtiene, se posee, “la gente anhela, busca y adquiere conocimiento” \cite[p.~126]{skinner1974sobre}. El aprendizaje se da cuando hay una modificación del comportamiento, generalmente observable, y que es generado a partir de estímulos y respuestas, mediante la práctica.

Las premisas pedagógicas que surgen de este enfoque son: 


\begin{itemize}
    
    \color{blue} \item \color{black} El docente entrega el conocimiento. 


    
     \color{blue} \item \color{black} El conocimiento se obtiene mediante el análisis de las partes para comprender la totalidad. 
     
    \color{blue} \item \color{black} El conocimiento es una verdad que debe ser adquirida para ser ejecutada. 
    \end{itemize}


\subsubsection*{Enfoque crítico}

Según este enfoque pedagógico, el conocimiento es un motor para la transformación social mediante la relación de la teoría con la práctica en contextos determinados. Su fin es la comprensión de la sociedad para mejorarla en pro de la calidad de vida. Difiere del cognitivismo y el conductismo, al plantear que el conocimiento se construye en y para un contexto determinado, mediante el proceso de interacción social entre la teoría, la práctica y la realidad. “Aprender significa generar cultura (...) generar investigación y reflexión de anuncio y denuncia, de acción y promoción, de pensamiento dialéctico en función de un objeto de estudio” \cite[p.~191]{ordonez2002pedagogia}. 

Las premisas pedagógicas de este enfoque son: 

\begin{itemize}
    
    \color{blue} \item \color{black} La flexibilidad o capacidad de adoptar formas diversas determinadas por las necesidades vitales y los problemas nacionales. 

     \color{blue} \item \color{black} La creatividad en la que “el estudiante deje de ser pasivo (...) y luchar por la participación, la recreación y la imaginación científica” \cite{Gonzalez2014}.
     
    \color{blue} \item \color{black} La dialogicidad mediante “una estructura pedagógica dialogal que implica horizontalidad, (...) crítica” \cite{Gonzalez2014}. 
    \end{itemize}
    
No es la intención hacer una descripción detallada sobre estos enfoques, pero es importante indicar que, a lo largo del proceso, se consideró necesaria la formación continua del docente de la carrera en aspectos didácticos y pedagógicos. Estos breves aspectos que han sido resaltados brindan una oportunidad de comprender mejor lo que se considera debe ser un estudiante de la carrera de Matemática, así como la manera de llevar a cabo la docencia. 

\subsection{Conceptualización de la población estudiantil y docente}

En este apartado se inicia con la conceptualización, tanto de la población estudiantil como del docente, para posteriormente tratar de entender la relación entre ambas poblaciones.

\subsubsection*{Concepción de la población estudiantil}

La población estudiantil de la carrera de Matemática de la Universidad de Costa Rica está constituida por todos los estudiantes que tengan la condición de estudiante activo o inactivo (que no está matriculado, aunque podría estar en el padrón estudiantil y por ello tener el derecho a la inscripción en cursos). Se concibe también una población extendida compuesta por egresados (quienes concluyeron el plan de estudios, pero tienen pendiente su trabajo final de graduación) y población graduada (quienes han obtenido su título de graduación en el grado académico de bachillerato). 

En uno de los talleres realizados, con la participación de docentes y de estudiantes avanzados de la carrera, se establecieron características que son necesarias que un estudiante presente, resumidas en la frase: “que sea responsable y disciplinado”. En relación con esta forma de concebir al estudiante, expresaron que “hay gente genial pero no tiene la disciplina de trabajo y conforme avanza la carrera esa genialidad por sí misma deja de ser suficiente”.

Otras características que se consideran relevantes son: que sea paciente, que practique el compañerismo, ya que es fundamental para el avance y el apoyo académico, por cuanto se necesita la cooperación entre los compañeros por encima de la competitividad. No se pretende un estudiante ideal porque no existe, pero se aspira a que sea alguien que no tenga un ego desmesurado y que sea persistente.

Existen dos realidades de estudiantes: por un lado, aquellos que son quejosos, dramáticos y que solo están interesados en tener un buen promedio; y, por otro, los esforzados y con deseos de superación. En los últimos años, se han incorporado diferentes actividades, entre las que se destacan los Coloquios y el Café Matemático, que brindan acompañamiento y mejores herramientas de trabajo al estudiante, lo cual le permite proyectarse como estudiante, y, asumir un rol más activo y comprometido con las exigencias que la carrera demanda. 

¿Cuál bagaje conceptual trae el estudiante al principio? En matemática es muy difícil diagnosticar esto porque vienen de realidades muy diversas. Lo primero que debe aprender no tiene que ver con los contenidos, es necesario que aprenda formas de organizar el pensamiento en matemática. 

\subsubsection*{Rol docente facilitador}

Según lo expresado por Rizzo \cite[p.~15]{rizzo2017ser}, “es muy importante lo que “sabe” un docente, así como la gestión de la clase: su “actitud”, las relaciones afectuosas y comunicativas que proponga, el logro de un “clima matemático” en el aula, para posibilitar y potenciar un aprendizaje significativo”. 

Dicha autora expresa cuál conocimiento debe poseer el docente de matemática para dar seguimiento al perfil de salida propuesto. Este conocimiento debe considerar lo siguiente: 

\begin{itemize}
    
    \color{blue} \item \color{black} El contexto, el cual implica reflexionar acerca de la dimensión social, política, cultural de la institución y de los alumnos, para una toma de decisiones del proceso de aprendizaje. Se reconoce que el docente es parte de los agentes de transformación en medio de una sociedad y de ahí el rol activo que conlleva estar atentos a las circunstancias actuales. 

     \color{blue} \item \color{black} Lo curricular sobre la matemática, que, en este caso, es de relevancia el perfil de salida determinado y su continua revisión al momento de planear un curso. 
     
    \color{blue} \item \color{black} Lo didáctico, en lo cual se consideran los espacios para la selección y uso de materiales didácticos como herramientas pedagógicas, así como que se logre dar respuesta a los propósitos declarados en la carrera. 
    
    \color{blue} \item \color{black} Sobre sí mismo, porque importa qué concepción de matemática tiene, qué le aporta a su forma de dar las clases, qué lo motiva y qué necesita mejorar, así como aquello que lo limita pero que puede ser una oportunidad de apoyarse en otros para alcanzar el fin deseado.
    \end{itemize}
    
Algunos docentes de la Escuela de Matemática de la Universidad de Costa Rica, en un taller realizado el 22 de mayo del 2019, comentaron que, durante mucho tiempo, se han dado prácticas como las siguientes: brindar listas de ejercicios, enseñar técnicas de exploración científica para realizar la investigación, proponer proyectos de investigación donde aprendan a escribir, comunicar, la forma de citar. Hay también exámenes orales y tareas que se exponen en clases. Para llevar a cabo todo esto, se incluye la siguiente lista de características que son necesarias y que fueron expresadas por dichos docentes: que tenga apertura, sepa escuchar, que pueda proyectar la voz, que sea ordenado, con buena letra, versátil, que busque diferentes formas de explicar cuando no se logra entender algo. Se reconoce que, con el paso del tiempo, se aprende que es bueno preguntar antes de borrar la pizarra, que se tenga pasión por lo que enseña, humildad académica, paciencia y, con mayor importancia, que se preparen las clases, porque el estudiante mide el conocimiento del profesor. 

En este rubro se mencionan valores como la paciencia, la perseverancia, la apertura, la cultura de trabajo y la autocrítica. Por otro lado, el trabajo colaborativo se menciona como una dinámica que se busca aplicar en los procesos de aprendizaje.  Por lo tanto, a partir de las tres conceptualizaciones anteriores, cada docente que desee innovar en sus aulas podría plantearse la siguiente premisa orientadora: asumir un enfoque pedagógico, según \cite{Gonzalez2014}, implica: 

\begin{itemize}
    
    \color{blue} \item \color{black} Ser consciente de las actividades pedagógicas que da más gusto realizar.

     \color{blue} \item \color{black} Determinar cuál es el enfoque pedagógico favorito, lo cual implica conocer algunos enfoques pedagógicos básicos para tener insumos en el momento de elegir.
     
    \color{blue} \item \color{black} Encontrar el enfoque pedagógico base del perfil profesional de la carrera en la cual se colabora.
    
    \color{blue} \item \color{black} Conocer, como punto de partida para innovar, las principales características didácticas de ese enfoque. 
    
    \color{blue} \item \color{black} Buscar la coherencia entre el mandato curricular y la postura pedagógica personal. 
    \end{itemize}

\subsubsection*{Relación docente-estudiante}

Considerando lo que se expresó en el taller del 22 de mayo del 2019, se destaca la empatía con los estudiantes como punto de partida, porque así se logra fomentar la participación para fortalecer la comunicación oral. La experiencia de aprendizaje se da tanto individual como grupalmente, y hay profesores que resaltan este último, otros no. Una clave de esto radica en las discusiones en grupo cuando ya se ha estudiado. Usualmente, la clase ha sido magistral, pero, el ideal de la clase es que los estudiantes lleguen estudiados y que se resuelvan ejercicios. 

Se reconoce que el esfuerzo grande debe ser en los primeros años, y, por este motivo, se procura la presencia de profesores convencidos de esta situación en los primeros años de carrera, donde se motive al estudiantado a que sean más cooperativos y menos egoístas, que tengan una cercanía entre ellos y que se mantenga la convicción de que muchos de los estudiantes serán futuros colegas. 

La carrera ha estado muy enfocada en resolver exámenes, entonces los estudiantes se adiestraron en eso. A largo plazo, se lograba desarrollar el pensamiento lógico, se adquiere madurez académica, profundidad analítica, un desarrollo del pensamiento abstracto. 

Si el profesor brinda apuntes del curso, bien claros, es una gran ayuda, porque, para hacer ejercicios, lo más importante es tener la materia, saber de cuáles fuentes se puede estudiar, hacer listas de ejercicios; pero, se menciona que Internet ha perjudicado un poco porque los estudiantes buscan cómo tener sus listas de ejercicios resueltas, no en asegurarse que van aprendiendo. 

Se destaca que el profesor debe distinguirse, sin necesidad de establecer una jerarquía específica. Por cómo se han dado los acontecimientos a lo largo de la historia, no existe posibilidad de que, en la formación inicial, un estudiante construya conocimiento desde cero. Por eso, el profesor guía ese acompañamiento teórico para que se lleguen a alcanzar los conceptos matemáticos. Por lo que se insiste en que el estudiante debe poseer pasión por la materia, que sea creativo, curioso, entre otras; sin embargo, en los primeros años, se ha dado el panorama que el estudiante casi ni participa, teme fallar, no soporta ser evidenciado ante los demás por no tener ese conocimiento que otros tal vez sí han adquirido previo a la carrera, por lo que se ha establecido una ruptura de la relación vertical entre el docente y el estudiante; esto provoca que, poco a poco, se vaya constituyendo una relación horizontal, con el apoyo de asistentes o de consejeros, con el fin de que el estudiante logre aceptar que, en la experiencia de aprender, en matemática, es necesario expresar las creencias que se poseen y poder, así, transformarlas con ayuda de los compañeros y del colectivo al cual decidió pertenecer. 

En esta relación, la postura de Hernández \cite{hernandez2017metodo} es la que se concibe dentro de las funciones del docente en su labor didáctica, para lo cual, se observa la necesidad de que, en el proceso de enseñanza y aprendizaje de la matemática, se consideren los siguientes componentes: 

\begin{itemize}
    
    \color{blue} \item \color{black} Motivación.

     \color{blue} \item \color{black} La orientación hacia el objetivo. 
     
    \color{blue} \item \color{black} El aseguramiento del nivel de partida. 
    
    \color{blue} \item \color{black} La elaboración del nuevo contenido.  
    
    \color{blue} \item \color{black} La fijación. 
    
    \color{blue} \item \color{black} El control y valoración del rendimiento.
    \end{itemize}

Estos componentes se incorporan para que el docente recuerde que la labor no es solo dentro del aula o en el momento de estar en comunicación con el estudiante. No obstante, esto debe procurarse en ambos sentidos: el estudiante también debe reflexionar acerca de su forma de aprender y qué debe mejorar. Por eso es que, en el perfil, se vislumbra a una persona que esté consciente de los conocimientos, habilidades y actitudes que le son propias de la carrera de Matemática. El docente es solo un guía que procura que dicho estudiante pueda comprender qué significa estar en Matemática.  


\subsection{Sobre la virtualidad}

El tema de la virtualidad, es decir, la posibilidad de ofrecer clases en línea (de forma sincrónica o asincrónica), se ha debatido ampliamente a lo largo del tiempo. Adquirió especial relevancia a partir de la pandemia de COVID-19 en 2020 y continúa vigente en la actualidad. La evidencia parece indicar que el rendimiento académico \textit{inmediato} (evaluado al final del curso) no difiere significativamente entre la modalidad virtual y la presencial \cite{hachey2022post}.

Las clases virtuales ofrecen ventajas tales como la posibilidad de grabarse, de modo que la persona estudiante puede revisarlas posteriormente y repetir aquellas partes que le resulten difíciles. También resultan convenientes para quienes encuentran complicado el traslado hasta la universidad, ya que ahorran tiempo y dinero al recibir la clase desde su lugar de residencia.

Sin embargo, también presentan problemas, como la falta de contacto directo con la persona docente y, en consecuencia, la pérdida de una discusión continua sobre el tema. Podría argumentarse que esta interacción también es posible en el ambiente virtual sincrónico, pero la experiencia observada por las personas docentes de matemáticas indica que rara vez ocurre. Esta pérdida de interacción reduce la profundidad del aprendizaje: aunque algunos estudios comparan las notas obtenidas en ambas modalidades y no encuentran diferencias significativas, lo cierto es que el estudiantado tiene una experiencia formativa más limitada al no darse esas interacciones directas. Por ejemplo, en un estudio realizado en Arabia Saudita \cite{Alshehri31122023}, se encontró que "students perceived technology as a supportive tool that could not replace the need for teacher–student interaction and engagement. In addition, online students perceived communication as an essential classroom element. Participants considered online assessment measures ineffective without useful feedback."

La diferencia es palpable cuando se hace una demostración. La experiencia de ver la construcción en vivo del elemento fundamental de la matemática pura y poder preguntar, señalar, opinar, ofrecer alternativas, proponer generalizaciones, etc. con retroalimentación inmediata de la persona docente, enriquece significativamente el proceso de aprendizaje en matemática. A esto podemos agregar las dificultades técnicas de interactuar por medio de la escritura en vivo con estudiantes para analizar un argumento.

En relación con la incorporación de componentes virtuales en la enseñanza de la matemática, resulta pertinente considerar que las características propias de la disciplina pueden plantear retos particulares para la educación a distancia. En este sentido, Yorkovsky y Levenberg \parencite{YorkovskyLevenberg2022}, al estudiar las experiencias de estudiantes de formación docente en ciencias y matemática, encontraron una clara preferencia por la enseñanza presencial frente a las modalidades a distancia, siendo la modalidad sincrónica la menos favorecida. Entre las principales dificultades señaladas se encuentran la comprensión de los contenidos, la interacción con la persona docente, los problemas técnicos y las limitaciones de los medios utilizados para ilustrar y representar conceptos matemáticos. Estos resultados sugieren que la incorporación de virtualidad en una carrera de Matemática Pura debe realizarse de manera reflexiva, atendiendo no solamente a la disponibilidad de recursos tecnológicos, sino también a las necesidades pedagógicas propias de una disciplina en la que la abstracción, el razonamiento simbólico, la discusión de argumentos y la interacción directa durante la resolución de problemas desempeñan un papel central.

Otra posible situación problemática asociada a las clases virtuales, en particular a las sincrónicas, es la posibilidad de una conexión deficiente o incluso desconexión debido a factores externos no controlables por la persona docente o estudiante. 

En cursos prácticos, donde el uso de la computadora en vivo es factible, podría ser que la modalidad virtual tenga más ventajas. 

En vista de lo anterior, se ha optado por un modelo basado mayoritariamente en clases presenciales. Se ha establecido que las modalidades virtuales (parcial o total) son apropiadas en algunas clases prácticas, con estilo de laboratorio computacional, donde puedan ofrecer ventajas. Se espera que esto contribuya a una formación más integral del estudiantado y tenga repercusiones positivas en el largo plazo.

\section{Cursos}

En esta sección se incluye la lista de cursos de la propuesta, organizada por año y, cuando corresponde, por énfasis. Cada entrada enlaza al programa recogido en la sección~\ref{sec:anexo-programas}.

\subsection{Cursos de primer año}

La siguiente tabla muestra los niveles de dominio y contenidos esperados para estos dos semestres.

\begin{table}[H]
\begin{center}
\begin{tabular}{|ccc|}
\hline
\multicolumn{3}{|c|}{\cellcolor[HTML]{00C0F3} \textbf{Niveles de dominio y contenidos por área curricular}} \\ \hline
\multicolumn{1}{|l|}{\cellcolor[HTML]{8ED8F8}\textbf{Área curricular}} & \multicolumn{1}{l|}{\cellcolor[HTML]{8ED8F8}\textbf{Dominio}} & \cellcolor[HTML]{8ED8F8}\textbf{Contenidos} \\ \hline
\multicolumn{1}{|c|}{Habilidades operacionales} & \multicolumn{1}{c|}{Alto} & Alto \\ \hline
\multicolumn{1}{|l|}{Pensamiento abstracto} & \multicolumn{1}{c|}{Medio} & Bajo \\ \hline
\multicolumn{1}{|l|}{Pensamiento algorítmico} & \multicolumn{1}{c|}{Bajo} & Bajo \\ \hline
\multicolumn{1}{|l|}{Formalismo} & \multicolumn{1}{c|}{Bajo} & Bajo \\ \hline
\multicolumn{1}{|c|}{Habilidades investigativas} & \multicolumn{1}{c|}{-} & Bajo \\ \hline
\end{tabular}
\end{center}
\caption{Niveles de dominio y contenidos para los cursos del primer año de la carrera.}
\end{table}

\begin{itemize}
    \item \hyperlink{MA0151}{MA-0151 Fundamentos de álgebra, trigonometría y geometría analítica}
    \item \hyperlink{MA0152}{MA-0152 Matemática Exploratoria}
    \item \hyperlink{MA0251}{MA-0251 Introducción al Cálculo en una variable}
    \item \hyperlink{MA0252}{MA-0252 Introducción a las Demostraciones}
\end{itemize}


\subsection{Cursos de segundo año}

\begin{table}[H]
\begin{center}
\begin{tabular}{|ccc|}
\hline
\multicolumn{3}{|c|}{\cellcolor[HTML]{00C0F3} \textbf{Niveles de dominio y contenidos por área curricular}} \\ \hline
\multicolumn{1}{|l|}{\cellcolor[HTML]{8ED8F8}\textbf{Área curricular}} & \multicolumn{1}{l|}{\cellcolor[HTML]{8ED8F8}\textbf{Dominio}} & \cellcolor[HTML]{8ED8F8}\textbf{Contenidos} \\ \hline
\multicolumn{1}{|l|}{Habilidades operacionales} & \multicolumn{1}{c|}{Alto} & Medio \\ \hline
\multicolumn{1}{|l|}{Pensamiento abstracto} & \multicolumn{1}{c|}{Medio} & Medio \\ \hline
\multicolumn{1}{|l|}{Pensamiento algorítmico} & \multicolumn{1}{c|}{Medio} & Medio \\ \hline
\multicolumn{1}{|l|}{Formalismo} & \multicolumn{1}{c|}{Medio} & Medio \\ \hline
\multicolumn{1}{|l|}{Habilidades investigativas} & \multicolumn{1}{c|}{-} & Bajo \\ \hline
\end{tabular}
\end{center}
\caption{Niveles de dominio y contenidos para los cursos del segundo año de la carrera.}
\end{table}

\begin{itemize}
    \item \hyperlink{MA0341}{MA-0341 Matemática Computacional}
    \item \hyperlink{MA0351}{MA-0351 Cálculo en varias variables}
    \item \hyperlink{MA0361}{MA-0361 Algebra Lineal I}
    \item \hyperlink{MA0451}{MA-0451 Principios de análisis I}
    \item \hyperlink{MA0461}{MA-0461 Algebra Lineal II}
\end{itemize}

\subsubsection*{Énfasis en matemática pura}
\begin{itemize}
    \item \hyperlink{MA0471}{MA-0471 Introducción a la Geometría Diferencial}
    \item \hyperlink{MA0496}{MA-0496 Teoría de Números}
\end{itemize}

\subsubsection*{Énfasis en matemática aplicada}
\begin{itemize}
    \item \hyperlink{CA0203}{CA-0203 Herramientas de ciencia de datos I}
    \item \hyperlink{CA0721}{CA-0721 Probabilidad}
\end{itemize}




\subsection{Cursos de tercer año}

La siguiente tabla resume los niveles de dominio y contenidos esperados para el tercer año, según las áreas curriculares definidas en este capítulo.

\begin{table}[H]
\begin{center}
\begin{tabular}{|ccc|}
\hline
\multicolumn{3}{|c|}{\cellcolor[HTML]{00C0F3} \textbf{Niveles de dominio y contenidos por área curricular}} \\ \hline
\multicolumn{1}{|l|}{\cellcolor[HTML]{8ED8F8}\textbf{Área curricular}} & \multicolumn{1}{l|}{\cellcolor[HTML]{8ED8F8}\textbf{Dominio}} & \cellcolor[HTML]{8ED8F8}\textbf{Contenidos} \\ \hline
\multicolumn{1}{|l|}{Habilidades operacionales} & \multicolumn{1}{c|}{Alto} & Bajo \\ \hline
\multicolumn{1}{|l|}{Pensamiento abstracto} & \multicolumn{1}{c|}{Alto} & Alto \\ \hline
\multicolumn{1}{|l|}{Pensamiento algorítmico} & \multicolumn{1}{c|}{Alto} & Alto \\ \hline
\multicolumn{1}{|l|}{Formalismo} & \multicolumn{1}{c|}{Alto} & Medio \\ \hline
\multicolumn{1}{|l|}{Habilidades investigativas} & \multicolumn{1}{c|}{Bajo} & Bajo \\ \hline
\end{tabular}
\end{center}
\caption{Niveles de dominio y contenidos para los cursos del tercer año de la carrera.}
\end{table}

\begin{itemize}
    \item \hyperlink{MA0515}{MA-0515 Ecuaciones Diferenciales}
    \item \hyperlink{MA0551}{MA-0551 Principios de Análisis II}
    \item \hyperlink{MA0635}{MA-0635 Topología General}
    \item \hyperlink{MA0641}{MA-0641 Análisis Numérico I}
\end{itemize}

\subsubsection*{Énfasis en matemática pura}
\begin{itemize}
    \item \hyperlink{MA0561}{MA-0561 Algebra Abstracta I}
    \item \hyperlink{MA0661}{MA-0661 Algebra Abstracta II}
\end{itemize}

\subsubsection*{Énfasis en matemática aplicada}
\begin{itemize}
    \item \hyperlink{CA0303}{CA-0303 Estadística I}
    \item \hyperlink{MA0647}{MA-0647 Modelación Matemática}
\end{itemize}

\subsection{Cursos de cuarto año}

La siguiente tabla resume los niveles de dominio y contenidos esperados para el cuarto año, según las áreas curriculares definidas en este capítulo.

\begin{table}[H]
\begin{center}
\begin{tabular}{|ccc|}
\hline
\multicolumn{3}{|c|}{\cellcolor[HTML]{00C0F3} \textbf{Niveles de dominio y contenidos por área curricular}} \\ \hline
\multicolumn{1}{|l|}{\cellcolor[HTML]{8ED8F8}\textbf{Área curricular}} & \multicolumn{1}{l|}{\cellcolor[HTML]{8ED8F8}\textbf{Dominio}} & \cellcolor[HTML]{8ED8F8}\textbf{Contenidos} \\ \hline
\multicolumn{1}{|l|}{Habilidades operacionales} & \multicolumn{1}{c|}{Alto} & Bajo \\ \hline
\multicolumn{1}{|l|}{Pensamiento abstracto} & \multicolumn{1}{c|}{Alto} & Alto \\ \hline
\multicolumn{1}{|l|}{Pensamiento algorítmico} & \multicolumn{1}{c|}{Alto} & Alto \\ \hline
\multicolumn{1}{|l|}{Formalismo} & \multicolumn{1}{c|}{Alto} & Alto \\ \hline
\multicolumn{1}{|l|}{Habilidades investigativas} & \multicolumn{1}{c|}{Medio} & Medio \\ \hline
\end{tabular}
\end{center}
\caption{Niveles de dominio y contenidos para los cursos del cuarto año de la carrera.}
\end{table}

\begin{itemize}
    \item \hyperlink{MA0701}{MA-0701 Seminario de Matemática}
    \item \hyperlink{MA0702}{MA-0702 Análisis Complejo}
    \item \hyperlink{MA0705}{MA-0705 Análisis Real I}
    \item \hyperlink{MA0780}{MA-0780 Comunicación en las ciencias}
\end{itemize}

\subsubsection*{Énfasis en matemática pura}
\begin{itemize}
    \item \hyperlink{MA0781Pura}{MA-0781 Práctica profesional en matemática pura}
\end{itemize}

\subsubsection*{Énfasis en matemática aplicada}
\begin{itemize}
    \item \hyperlink{CA0411AnalisisDatos}{CA-0411 Análisis de Datos I}
    \item \hyperlink{MA0781Aplicada}{MA-0781 Práctica profesional en matemática aplicada}
\end{itemize}

\subsection{Cursos de licenciatura}

Las áreas curriculares del capítulo de estructura curricular contemplan explícitamente los años I a IV; para la licenciatura se proyecta la culminación de esa progresión, con énfasis reforzado en habilidades investigativas mediante los seminarios de estudios dirigidos.

\begin{table}[H]
\begin{center}
\begin{tabular}{|ccc|}
\hline
\multicolumn{3}{|c|}{\cellcolor[HTML]{00C0F3} \textbf{Niveles de dominio y contenidos por área curricular}} \\ \hline
\multicolumn{1}{|l|}{\cellcolor[HTML]{8ED8F8}\textbf{Área curricular}} & \multicolumn{1}{l|}{\cellcolor[HTML]{8ED8F8}\textbf{Dominio}} & \cellcolor[HTML]{8ED8F8}\textbf{Contenidos} \\ \hline
\multicolumn{1}{|l|}{Habilidades operacionales} & \multicolumn{1}{c|}{Alto} & Bajo \\ \hline
\multicolumn{1}{|l|}{Pensamiento abstracto} & \multicolumn{1}{c|}{Alto} & Alto \\ \hline
\multicolumn{1}{|l|}{Pensamiento algorítmico} & \multicolumn{1}{c|}{Alto} & Alto \\ \hline
\multicolumn{1}{|l|}{Formalismo} & \multicolumn{1}{c|}{Alto} & Alto \\ \hline
\multicolumn{1}{|l|}{Habilidades investigativas} & \multicolumn{1}{c|}{Alto} & Alto \\ \hline
\end{tabular}
\end{center}
\caption{Niveles de dominio y contenidos proyectados para los cursos de licenciatura.}
\end{table}

\begin{itemize}
    \item \hyperlink{MA0781SED}{MA-0781 Seminario de estudios dirigidos en matemática}
    \item \hyperlink{MA0881}{MA-0881 Seminario de estudios dirigidos en matemática}
\end{itemize}

\subsection{Cursos optativos}
\subsubsection*{Cursos optativos de núcleo en matemática pura OPT-001}

\begin{itemize}
    \item {MA-0788 Análisis Real II}
    \item \hyperlink{MA0506}{MA-0506 Teoría algebraica de números}
    \item {MA-0506 Teoría analítica de números}
    \item \hyperlink{MA0512}{MA-0512 Teoría de conjuntos}
    \item \hyperlink{MA0525}{MA-0525 Combinatoria}
    \item \hyperlink{MA0528}{MA-0528 Sistemas dinámicos}
    \item \hyperlink{MA0703}{MA-0703 Integración}
    \item \hyperlink{MA0707}{MA-0707 Geometría Diferencial}
    \item \hyperlink{MA0709}{MA-0709 Geometría algebraica}
    \item \hyperlink{MA0711}{MA-0711 Lógica}
    \item \hyperlink{MA0714}{MA-0714 Optimización}
    \item \hyperlink{MA0755}{MA-0755 Ecuaciones diferenciales parciales}
    \item \hyperlink{MA0804}{MA-0804 Topología algebraica}
    \item \hyperlink{MA0806}{MA-0806 Análisis funcional}
    \item \hyperlink{MA0815}{MA-0815 Análisis armónico}
    \item \hyperlink{MA0820}{MA-0820 Teoría de modelos}
    \item \hyperlink{MA0840}{MA-0840 Probabilidad}
    \item \hyperlink{MA0860}{MA-0860 Teoría de módulos}
    \item \hyperlink{MA0889}{MA-0889 Álgebra conmutativa}
    \item \hyperlink{MA0920}{MA-0920 Ecuaciones en derivadas parciales numéricas}
\end{itemize}

\subsubsection*{Cursos optativos de núcleo en matemática pura OPT-002}
\begin{itemize}
    \item {CA-304 Herramientas de ciencia de datos II}
    \item {CA-0411 Análisis de Datos II}
    \item {CA-0403 Estadística actuarial II}
    \item {CA-0512 Modelos Lineales}
    \item {CA-0612 Series de Tiempo}
    \item \hyperlink{MA0406}{MA-0406 Introducción a la optimización}
    \item {MA-0477 Ecuaciones Diferenciales Aplicadas}
    \item {MA-0722 Estadística Matemática I}
\end{itemize}

\subsubsection*{Cursos optativos en general OPT-003}


\begin{itemize}
    \item {MA-0788 Análisis Real II}
    \item \hyperlink{MA0506}{MA-0506 Teoría algebraica de números}
    \item {MA-0506 Teoría analítica de números}
    \item \hyperlink{MA0512}{MA-0512 Teoría de conjuntos}
    \item \hyperlink{MA0525}{MA-0525 Combinatoria}
    \item \hyperlink{MA0528}{MA-0528 Sistemas dinámicos}
    \item \hyperlink{MA0703}{MA-0703 Integración}
    \item \hyperlink{MA0707}{MA-0707 Geometría Diferencial}
    \item \hyperlink{MA0709}{MA-0709 Geometría algebraica}
    \item \hyperlink{MA0711}{MA-0711 Lógica}
    \item \hyperlink{MA0714}{MA-0714 Optimización}
    \item \hyperlink{MA0755}{MA-0755 Ecuaciones diferenciales parciales}
    \item \hyperlink{MA0804}{MA-0804 Topología algebraica}
    \item \hyperlink{MA0806}{MA-0806 Análisis funcional}
    \item \hyperlink{MA0815}{MA-0815 Análisis armónico}
    \item \hyperlink{MA0820}{MA-0820 Teoría de modelos}
    \item \hyperlink{MA0840}{MA-0840 Probabilidad}
    \item \hyperlink{MA0860}{MA-0860 Teoría de módulos}
    \item \hyperlink{MA0889}{MA-0889 Álgebra conmutativa}
    \item \hyperlink{MA0920}{MA-0920 Ecuaciones en derivadas parciales numéricas}
    \item {CA-304 Herramientas de ciencia de datos II}
    \item {CA-0411 Análisis de Datos II}
    \item {CA-0403 Estadística actuarial II}
    \item {CA-0512 Modelos Lineales}
    \item {CA-0612 Series de Tiempo}
    \item \hyperlink{MA0406}{MA-0406 Introducción a la optimización}
    \item {MA-0477 Ecuaciones Diferenciales Aplicadas}
    \item {MA-0722 Estadística Matemática I}    
    \item \hyperlink{MA0609}{MA-0609 Teoría analítica de números}
    \item \hyperlink{MA0790}{MA-0790 Tópicos de análisis}
    \item \hyperlink{MA0791}{MA-0791 Tópicos de Probabilidad}
    \item \hyperlink{MA0792}{MA-0792 Tópicos de Ecuaciones diferenciales}
    \item \hyperlink{MA0793}{MA-0793 Tópicos de Geometría}
    \item \hyperlink{MA0794}{MA-0794 Tópicos de Modelación}
    \item \hyperlink{MA0796}{MA-0796 Tópicos de Álgebra}
    \item \hyperlink{MA0797}{MA-0797 Tópicos de Matemática discreta}
    \item \hyperlink{MA0830}{MA-0830 Tópicos de Teoría de Números}
    \item \hyperlink{MA0831}{MA-0831 Tópicos de Lógica}
    \item \hyperlink{MA0832}{MA-0832 Tópicos en Topología}
    \item \hyperlink{MA0918}{MA-0918 Procesos estocásticos}
    \item \hyperlink{MA0794}{MA-0794 Tópicos de Modelación}
    \item \hyperlink{MA0795}{MA-0795 Tópicos de Computación científica}
    \item \hyperlink{MA0798}{MA-0798 Tópicos de Análisis de datos}
    \item \hyperlink{MA0799}{MA-0799 Tópicos de Aprendizaje automático}
    \item {MA-0759 Tópicos de Matemática Financiera}
    \item \hyperlink{MA0917}{MA-0917 Estadística II}
    \item \hyperlink{MA0918}{MA-0918 Procesos estocásticos}
    \item {CA-917 Estadística matemática II}
    \item \hyperlink{MA0795}{MA-0795 Tópicos de Computación científica}
\end{itemize}


